\section{Multiplicity and exactness}
\label{sec:combination}
Fix $b\in(0,2)$ and a compact $J\Subset(0,2)$ containing $b$.
The upper assertion supplies common $K,V,\rho_{\rm H}$ for $J$.
The two orbits constructed in Section~\ref{sec:lower-main}
belong to the same field, are distinct and isolated, and each
converges to the whole graphic. For all sufficiently large $n$,
both are contained in $K$, both satisfy the same Hausdorff
condition, and $\mu_n\in V$. Hence
\[
 1\le m(L_n^-),\qquad 1\le m(L_n^+),\qquad
 m(L_n^-)+m(L_n^+)\le2 .
\]
Both actual multiplicities are one, and there is no third isolated
whole-graphic cycle in the counted neighborhood.

Let $a_n^-<a_n^+$ be the two roots in
\eqref{eq:two-return-roots}, and write $\Pi_n=D_{p_n}\circ H_{p_n}$
for their actual minimal return. On its returning interval,
$D_{p_n}'>0$ and $H_{p_n}'>0$, so $\Pi_n'>0$.
Every zero gives an isolated cycle by the analytic-section argument
of Section~\ref{sec:lower-main}. The image estimates on the full
source rectangle place each such cycle in the same counted
neighborhood for all sufficiently large $n$. A further zero would
therefore contradict the total multiplicity bound. Hence these are
the only two zeros on the returning interval. Conjugacy to the
physical analytic section preserves the derivative at a fixed
point, so multiplicity one gives $\Pi_n'(a_n^\pm)-1\ne0$.
The negative--positive--negative displacement signs in
\eqref{eq:actual-lower-return} now imply
\[
 (\Pi_n-\mathrm{id})'(a_n^-)>0,\qquad
 (\Pi_n-\mathrm{id})'(a_n^+)<0.
\]
Consequently
\begin{equation}\label{eq:realizing-multipliers}
 \Pi_n'(a_n^-)>1,\qquad 0<\Pi_n'(a_n^+)<1.
\end{equation}
These are the characteristic multipliers of the planar positive-time
return. They are preserved by changes of transverse section or
section orientation. Thus $L_n^-$ is hyperbolic repelling and $L_n^+$
is hyperbolic attracting, with the labels ordered by the coordinate $a$.

The realizing sequence gives $C_{\rm H}(b)\ge2$ at each fixed
$b$. Assertion (i) gives cyclicity at most two at that same $b$,
and uniformly on compact modulus intervals. This completes
Theorem~\ref{thm:main}. The realizing sequence is chosen separately
for each fixed $b\in(0,2)$.
